% s5: who carries the note (points 1-3 of the theory). ASCII only.
\section*{5. Who carries the note}

The Hamiltonian commutes with the room's circle of symmetry, so it never mixes charges. Helike has charge $\pm2$ (the
$m=\pm2$ part of a spin-2, tensor-like object); her competitor has charge 0 (scalar-like).

\begin{claim}{PROVED}{Different charges cross exactly}\label{b:cross}
Helike and the competitor lie in different invariant sectors of the Berger circle, so their levels cross exactly and the
carrier is well defined on each side of a crossing.
\evidence{The circle commutes with the Hamiltonian at every $s$ and $g$; analytic perturbation theory per
sector~\cite{Kato}.}
\lines{17}
\end{claim}

\spine{2} Topology decides who is lightest. With the sources off, the potential of the room is the squared slope of
one function, the Chern--Simons number: Witten's form of a Schr\"odinger operator~\cite{Witten82}, whose low levels sit
at the critical points of that function, sorted by their number of downhill directions.

\begin{claim}{PROVED}{The landscape is the slope of Chern--Simons; the competitor is the topological direction}\label{b:susy}
Let $\omega=(\sqrt t,t^{-1/2},t^{-1/2})$, $S=\sum_k\omega_k$ and $W(q)=\tfrac12\sum_k\omega_k\|q_k\|^2-g\det q$. Then
$V_0=\tfrac12|\nabla W|^2$, $W=CS/(2\sqrt t\,g^2)$ on homogeneous fields, and $H=\tfrac12\sum_ia_i^*a_i+\tfrac32S$
exactly, $a_i=\partial_i+\partial_iW$; so the vacuum energy is $\tfrac32S$ up to exponentially small terms. Zero field is
the minimum of $W$; $q_{\rm MC}$ is a saddle of $W$ with one downhill direction, and that direction is the competitor,
pointing back along the path toward the other winding vacuum. Its frequency
$\omega_C=\tfrac12(\sqrt{t+8}-\sqrt t)$ ($0.7923$ at $t=3$) lies below Helike's $\omega_H=2/\sqrt t$ ($1.1547$) and every
other mode, for every $t>0$.
\evidence{$\partial\det q/\partial q_{kc}=\mathrm{cof}(q)_{kc}$ and exact block diagonalization of the Hessian, confirmed
by the referee pass~\cite{ClosedSet,Witten82}. Proved for the undressed reduced operator at $s=1$.}
\lines{8, 18}
\end{claim}

If I could put only one claim first it would be this: the cheapest way to excite the unsourced room is to lean toward
changing its topology. If it survives into the full theory, it is why the scalar glueball is lightest, which simulations
see~\cite{MorningstarPeardon99} and nobody has explained.

\spine{3} Tensor and scalar trade places at a calculable point. Turn the sources on ($s>1$): Helike's frequency
$\omega_H(s)^2=2s(1+2/t-s)$ falls toward the competitor's, and the order reverses.

\begin{claim}{PROVED}{The harmonic switch}\label{b:switch0}
At harmonic order the tensor (charge $\pm2$) and scalar (charge 0) carriers trade places at
$s_0^*(t)=(t+2)(\sqrt{t+8}-\sqrt t)/(4\sqrt t)$, which is $(5\sqrt{33}-15)/12=1.1436$ at $t=3$ and lies inside Theorem
A's window exactly when $t>1$.
\evidence{Closed-form root of the charge-0 secular equation~\cite{ClosedSet}.}
\lines{18}
\end{claim}

\begin{claim}{CONJECTURE}{The carrier changes only at the switch}\label{b:switch}
In the reduced model, for $s>s^*$ Helike carries the note and for $s<s^*$ the competitor does; at $s=1$ it is the
competitor. The registered switches $s^*=1.1674$, $1.1809$, $1.1931$ at $g^2=0.03$, $0.05$, $0.07$ rise from $s_0^*$ as
the coupling grows.
\evidence{The rise matches the anharmonic prediction at leading order; the registered values do not include the fast
dressing of \S6~\cite{ClosedSet}.}
\lines{18}
\end{claim}

\begin{claim}{CONDITIONAL}{A switch exists}\label{b:ivt}
At weak coupling ($g^2\lesssim0.0075$, $t=3$), given H8 at $s=1$ and Helike lowest at the controls (Box~\ref{b:order},
on H1), exact crossing and continuity give a switch. At the registered couplings the controls at $s=1.15$ already have the
competitor lowest, so existence there needs one more control at some $s_1\in(s^*,1.25)$ where Helike is lowest.
\evidence{Deduction from lines 12, 17, 19 and H8; the bracketing gap was found by the referee pass~\cite{ClosedSet}.}
\lines{18, 12}
\end{claim}

The sources are scaffolding; the physical room has none, and the bridge back is short.

\begin{claim}{PROVED}{The source bridge $s\to1^+$ (by deduction, reduced model)}\label{b:bridge}
Ordered eigenvalues and spectral projections of the reduced operator are continuous on $[1,s_1]$ for any $s_1$ in the
window.
\evidence{A type-(A) analytic family~\cite{Kato}. Only ordered eigenvalues are continuous through a charge-0 avoided
crossing near $s=1+0.95\,g^2$; this is not a bridge in the full theory~\cite{ClosedSet}.}
\lines{19}
\end{claim}

\textbf{The honest catch.} All of this is about the undressed reduced operator. In full Yang--Mills $q_{\rm MC}$ is a
gauge copy of $q=0$, so the two wells are one point; near $q=0$ the charge-0 and charge-$\pm2$ two-quantum states tie at
$2/\sqrt t$, and the tie is broken at order $g^2$ by a kinetic term of the fast modes (\S6).

\begin{claim}{OPEN}{H8 in the physical theory}\label{b:h8}
At $s=1$ the lowest excitation has charge 0. True for the undressed reduced operator (margin $0.3624$ at $t=3$); in the
physical theory decided only at $O(g^2)$, by the kinetic correction of Box~\ref{b:kinetic}.
\evidence{The register's correction to H8~\cite{ClosedSet}; the next step is a pair of one-sided bounds with constants
summing below $5.18$ at $g^2\le0.07$.}
\lines{18, 16}
\end{claim}

\textbf{Theta and tunnelling.} The degree-one gauge map shifts the Chern--Simons number by one, so the vacuum angle
$\theta$ is a Bloch phase~\cite{JackiwRebbi76,CallanDashenGross76}, and tunnelling between windings is set by the Agmon
distance~\cite{Agmon82,Simon83}.

\begin{claim}{PROVED}{The tunnelling length is pinned by topology}\label{b:agmon}
In variables $q=p/g$ the Agmon distance from zero field to the winding-one orbit satisfies
$1/(2\sqrt t)\le d\le\sqrt{(t+2)(2t+1)}/(6t)$ ($0.2887\le d\le0.3287$ at $t=3$), so tunnelling moves levels by
$O(e^{-2d/g^2})$. The lower bound is Bogomolny's: $\int|\nabla W|\,|dp|\ge\Delta W$.
\evidence{The analogue of the instanton bound $8\pi^2|\Delta CS|/g^2$~\cite{BPST75,tHooft76}; the upper bound is the
straight path~\cite{ClosedSet}.}
\lines{18, 19}
\end{claim}

I conjecture that theta barely moves the note, $|\Delta(\theta)-\Delta(0)|\le C\,e^{-8\pi^2/g^2}$ by tight binding over
the chain of vacua~\cite{tHooft76}; nothing in the theory uses it at $\theta=0$.

If the switch rises with the coupling, and a larger coupling is a larger room, the scalar's territory grows with the room.

\begin{claim}{CONJECTURE}{The small-room order is the glueball order}\label{b:glueball}
$s^*$ keeps rising along the growing room, and the small-room order, charge 0 below charge 2, becomes
$m(0^{++})<m(2^{++})$ in infinite volume.
\evidence{Lattice order~\cite{MorningstarPeardon99}. The rise is first order only, and continuing circle charges to
$J^{PC}$ needs an argument I do not have. I quote no mass ratios.}
\lines{18, 20}
\end{claim}

The carriers are named, topology orders them, and they trade places at a point I can write down, all read from nine
slow numbers. What of the infinitely many fast modes left out?
